% problem-000211 4 磷同素异形体晶体学
\section*{4. 磷同素异形体晶体学}
$\mathrm { Z n \mathrm { - } I _ { 2 } }$ 电池是⼀种⽔系⼆次电池，作为潜在的储能系统，近年来重新受到关注。经典的 $\mathrm { Z n \mathrm { - } I _ { 2 } }$ 电池正极为单质 $\mathrm { I } _ { 2 } ,$ ，负极为⾦属 $Z \mathrm { n }$ ，以 $\mathrm { Z n S O _ { 4 } }$ 溶液为电解质。研究发现，当负极区的电解质仍为 $\mathrm { Z n S O _ { 4 } }$ ，但正极区电解质变为 $\mathrm { C u S O _ { 4 } }$ 或 $\mathrm { C u C l _ { 2 } }$ 时，可调变电池正极的反应，对应的标准电极电势和放电容量均有显著变化，提⾼了该电池的储能密度。以I 为正极，⼏组代表性的充放电实验数据总结于表1中。⼀些氧化还原电对的标准电极电势和⽔合离⼦的标准Gibbs⽣成⾃由能变 $\Delta _ { \mathrm { f } } G _ { \mathrm { m } } ^ { \ominus }$ ⻅表2。

表 1. $\mathrm { I } _ { 2 }$ 电极在不同的正极区电解质中充放电实验结果

\textit{（原卷此处含表格/仪器试剂表，详见 PDF 或 MD 源文件。）}

表2. 可能⽤到的⼀些热⼒学数据

\textit{（原卷此处含表格/仪器试剂表，详见 PDF 或 MD 源文件。）}

\subsection*{4-1}
写出实验1\~4对应的正极区电极反应⽅程式。

\subsection*{4-2}
实验3放电⽣成⼀种含碘难溶电解质A，实验4充电⽣成⼀种含碘化合物 $\mathbf { B } _ { \circ }$

\subsection*{4-2}
-1 计算A的溶度积常数；

\subsection*{4-2}
-2 计算B的标准Gibbs⽣成⾃由能变。

\subsection*{4-3}
假设电解质中的离⼦浓度均恒为标态，计算下列两种 $\mathrm { Z n \mathrm { - } I _ { 2 } }$ 电池的储能密度（基于正负极的总质量，结果以Wh/kg表⽰）。

\subsection*{4-3}
-1 正负极均以 $\mathrm { Z n S O _ { 4 } }$ 为电解质；

\subsection*{4-3}
-2 负极以 $\mathrm { Z n S O _ { 4 } }$ 电解质，正极以 $\mathrm { C u C l _ { 2 } }$ 为电解质。

\subsection*{4-4 某 $\mathrm { Z n - I _ { 2 } }$ 电池可表⽰为：(− $\mathrm { ) Z n ( s ) \mid Z n S O _ { 4 } ( 1 . 0 M ) \mid \mid C u C l _ { 2 } ( 1 . 0 M ) \mid \mid I _ { 2 } ( s ) ( + ) }$}

电池正极区和负极区电解质溶液量均为1.0 mL，Zn和 $\mathrm { I } _ { 2 }$ 的量均为 2.0 mmol，以200 mA 恒电流放电 482$\mathsf { s } \varlimsup$ ，计算电池的电动势（计算时 $\mathrm { I } _ { 2 \mathrm { \Delta } }$ 、A和B均可以作为不溶物处理）。
